\section{Prove the accumulator width before choosing an instruction}
One signed halfword product can equal
\((-32768)(-32768)=1073741824=2^{30}\). That fits a signed word. Two such products
sum to \(2^{31}\), which does not. Eight sum to \(2^{33}\). The most negative
possible eight-product sum is \(-8589672448\), obtained with opposite-sign extrema.
The positive endpoint alone requires thirty-five signed bits; a signed thirty-four-bit
integer stops at \(2^{33}-1\).

Use an exact 64-bit accumulator. In the scalar path, sign-extend both halfwords,
form each exact product and add with \texttt{daddu}. Reset the accumulator for each
output. The native doubleword operations keep all relevant bits. Their wraparound
semantics do not lose information here because the proven complete and partial sums
fit a signed doubleword.

A signed halfword dot product widened to words still has the two-product problem.
For two minimum input samples and two minimum coefficients, a word result can wrap
from positive \(2^{31}\) to negative \(-2^{31}\). Copying that already-overflowed
word into a 64-bit register cannot recover the lost distinction. Widen before the
partial sum is formed.

\begin{table}[htbp]
  \centering
  \caption{Positive boundary sums from minimum signed sample and coefficient codes.}
  \tagpdfsetup{table/header-rows={1}}
  \begin{tabular}{rrl}
    \toprule
    Products & Exact sum  & Signed storage sufficient for this value \\
    \midrule
    1        & 1073741824 & 32 bits                                  \\
    2        & 2147483648 & 33 bits                                  \\
    4        & 4294967296 & 34 bits                                  \\
    8        & 8589934592 & 35 bits                                  \\
    \bottomrule
  \end{tabular}
\end{table}
